\section{Struct (Conjuntos Estruturados)}
\label{sec:const-struct}

A categoria \textbf{Struct} permite modelar sistemas dinâmicos discretos onde cada elemento possui um valor associado e um ponteiro para um sucessor (endomapa). As construções universais nesta categoria garantem a preservação da coerência entre os dados armazenados e a estrutura de transição.

\subsection{Fundamentação Teórica}

\begin{itemize}
    \item \textbf{Produto:} 
    O produto de $(d_1, \varphi_1, n_1)$ e $(d_2, \varphi_2, n_2)$ é o triplo $(d, \varphi, n_1 \cdot n_2)$, onde os elementos são indexados por pares $(i,j)$. A estrutura e os dados são definidos por:
    \[
        \varphi(i,j) = (\varphi_1(i), \varphi_2(j)), \qquad d(i,j) = (d_1(i), d_2(j)).
    \]

    \item \textbf{Coproduto:} O coproduto $S_1 \sqcup S_2$ é a união disjunta dos conjuntos de elementos, mantendo as funções $\varphi$ e $d$ originais em cada componente. O número total de elementos é $n_1 + n_2$.

    \item \textbf{Pullback:} Para um diagrama $S_1 \xrightarrow{f} S_3 \xleftarrow{g} S_2$, o pullback é o subobjeto do produto $S_1 \times S_2$ constituído pelos pares $(i,j)$ tais que $f(i) = g(j)$.

    \item \textbf{Pushout:} O pushout de $S_1 \xleftarrow{f} S_0 \xrightarrow{g} S_2$ é o coproduto $S_1 \sqcup S_2$ submetido à menor relação de equivalência que identifica $f(k) \sim g(k)$ para todo $k \in S_0$, estendendo-se para manter a compatibilidade com os endomapas.

    \item \textbf{Equalizador:} O equalizador de $f, g: S_1 \rightrightarrows S_2$ é o subobjeto de $S_1$ formado pelos elementos $i$ tais que $f(i) = g(i)$. Por definição de morfismo, este conjunto é automaticamente fechado para a aplicação de $\varphi_1$.

    \item \textbf{Coequalizador:} O coequalizador de $f, g: S_1 \rightrightarrows S_2$ é o quociente de $S_2$ pela menor congruência compatível com $\varphi_2$ que identifica $f(i) \sim g(i)$ para todo $i \in S_1$.
\end{itemize}

\subsection{Implementação Computacional}

Abaixo apresenta-se a implementação em MATLAB, respeitando as estruturas e algoritmos de união-busca presentes no script original.

\begin{lstlisting}[language=Matlab, caption={Produto e Coproduto em Struct}]
function P = struct_product(S1, S2)
    P.n   = S1.n * S2.n;
    P.phi = zeros(1, P.n);
    P.d   = zeros(1, P.n, 'like', 1i);
    for i = 1:S1.n
        for j = 1:S2.n
            k = (i-1)*S2.n + j;
            % phi(i,j) = (phi1(i), phi2(j))
            P.phi(k) = (S1.phi(i)-1)*S2.n + S2.phi(j);
            P.d(k)   = S1.d(i) + S2.d(j); 
        end
    end
end

function C = struct_coproduct(S1, S2)
    C.n   = S1.n + S2.n;
    C.phi = [S1.phi, S2.phi + S1.n];
    C.d   = [S1.d, S2.d];
end
\end{lstlisting}

\begin{lstlisting}[language=Matlab, caption={Equalizador e Pullback em Struct}]
function [E, inc] = struct_equalizer(S1, f, g)
    E_idx = find(f == g);
    E.n   = numel(E_idx);
    inc   = E_idx;
    E.d   = S1.d(E_idx);
    E.phi = zeros(1, E.n);
    for i = 1:E.n
        target = S1.phi(E_idx(i));
        [~, loc] = ismember(target, E_idx);
        E.phi(i) = loc;
    end
end

function [PB, p1, p2] = struct_pullback(S1, S2, f, g)
    % Baseado na construcao de subobjeto do produto
    pairs = [];
    for i = 1:S1.n
        for j = 1:S2.n
            if f(i) == g(j)
                pairs = [pairs; i, j];
            end
        end
    end
    PB.n = size(pairs, 1);
    p1 = pairs(:,1)'; p2 = pairs(:,2)';
    PB.d = S1.d(p1); % d1 = d2 o f no pullback
    PB.phi = zeros(1, PB.n);
    for k = 1:PB.n
        t1 = S1.phi(p1(k)); t2 = S2.phi(p2(k));
        [~, PB.phi(k)] = ismember([t1, t2], pairs, 'rows');
    end
end
\end{lstlisting}

\begin{lstlisting}[language=Matlab, caption={Coequalizador e Pushout em Struct}]
function [CE, quot] = struct_coequalizer(S1, S2, f, g)
    parent = 1:S2.n;
    function r = find_root(x)
        while parent(x) ~= x, x = parent(x); end
        r = x;
    end
    % Identificar f(i) ~ g(i)
    for i = 1:S1.n
        a = find_root(f(i)); b = find_root(g(i));
        if a ~= b, parent(a) = b; end
    end
    % Fechar para congruencia (phi)
    changed = true;
    while changed
        changed = false;
        for x = 1:S2.n
            rx = find_root(x);
            pa = find_root(S2.phi(x));
            pb = find_root(S2.phi(rx));
            if pa ~= pb
                parent(pa) = pb; changed = true;
            end
        end
    end
    [roots, ~, quot] = unique(arrayfun(@find_root, 1:S2.n));
    CE.n = numel(roots);
    CE.phi = quot(S2.phi(roots));
    CE.d = S2.d(roots);
end

function [PO, i1, i2] = struct_pushout(S0, S1, S2, f, g)
    C = struct_coproduct(S1, S2);
    f_c = f; % f em S1
    g_c = g + S1.n; % g em S2 deslocado
    [PO, quot] = struct_coequalizer(S0, C, f_c, g_c);
    i1 = quot(1:S1.n);
    i2 = quot(S1.n+1:end);
end
\end{lstlisting}